\documentclass[11pt,a4paper]{article}
\usepackage[utf8]{inputenc}
\usepackage[T1]{fontenc}
\usepackage[french]{babel}
\usepackage[margin=2cm]{geometry}
\usepackage{graphicx}
\usepackage{hyperref}
\usepackage{enumitem}
\usepackage{xcolor}
\usepackage{titlesec}

\setlist{nosep,leftmargin=1.2em}
\titlespacing*{\section}{0pt}{0.6em}{0.3em}
\titlespacing*{\subsection}{0pt}{0.4em}{0.2em}
\titleformat{\section}{\normalfont\bfseries\large}{\thesection}{0.5em}{}
\titleformat{\subsection}{\normalfont\bfseries}{\thesubsection}{0.5em}{}
\setlength{\parskip}{0.3em}
\setlength{\parindent}{0pt}

\title{\vspace{-1.5cm}Rapport individuel --- Sprint 2}
\author{Louis-Marie Simonneaux (\texttt{lsx4897}) --- Projet Long SHDL}
\date{10 mai 2026}

\begin{document}
\maketitle
\vspace{-1.5em}

\section*{Contexte g\'en\'eral}

Le projet vise \`a d\'evelopper un simulateur de circuits logiques d\'ecrits en SHDL, avec interface JavaFX. L'\'equipe est r\'epartie sur six axes (\'editeur, parser, automates, simulateur, IHM, int\'egration). Mon axe est la \textbf{repr\'esentation graphique du circuit} : c\^ot\'e utilisateur, la fen\^etre o\`u l'on bascule les entr\'ees et observe les sorties.

L'it\'eration~1 avait livr\'e un prototype JavaFX autonome (deux rang\'ees de \texttt{Toggle\-Button} ronds et carr\'es) qui d\'emontrait la ma\^itrise des conteneurs JavaFX et du style CSS, mais qui \textbf{n'\'etait reli\'e \`a aucun moteur de simulation}. Le bilan de sprint avait identifi\'e ce manque comme le point bloquant pour rejoindre le MVP bout-en-bout.

Le sprint~2 se d\'ecline donc en deux it\'erations~: (i) brancher la fen\^etre sur le simulateur r\'eel et amorcer la notion de module composite ; (ii) absorber la refonte de l'API du simulateur entreprise par l'\'equipe et int\'egrer mes fichiers d'interface dans le module commun.

\section{It\'eration 2 --- branchement au simulateur}

\subsection*{Objectif}

Brancher la fen\^etre prototype, jusque-l\`a purement cosm\'etique, sur l'API \texttt{Simulateur~/ Composant~/ Lien} fournie par l'axe simulateur, afin que les changements d'\'etat des entr\'ees se propagent \`a un vrai circuit logique.

\subsection*{Activit\'es r\'ealis\'ees}

\begin{itemize}
  \item \textbf{Refactoring de \texttt{FenetreSimulation}} (commit \texttt{62c4bbc}, +81/-64 lignes) : r\'eorganisation des m\'ethodes de construction de l'IHM, extraction de la logique de style CSS dans des constantes nomm\'ees (\texttt{StyleUpCarre}, \texttt{StyleDownRond}, \texttt{StyleUnknownCarre}), simplification du couplage entre les rang\'ees de \texttt{ToggleButton} et leurs handlers d'\'ev\'enement.
  \item \textbf{Branchement au simulateur} (commit \texttt{6c9bfe4}, 6~fichiers, +131/-41) : la fen\^etre instancie d\'esormais un \texttt{Simulateur(NB\_ENTREES, NB\_SORTIES)} et y ajoute un circuit de d\'emonstration constitu\'e de 4~portes \texttt{And} c\^abl\'ees aux \texttt{Lien} d'entr\'ee et de sortie via \texttt{sim.ajouterComposant(...)}. Les clics sur les carr\'es d'entr\'ee modifient les \'etats UP/DW/ND, le simulateur les propage et les ronds de sortie sont mis \`a jour en cons\'equence.
  \item \textbf{D\'ebut de la classe \texttt{Module}} (\texttt{tests projet long/simulateur/Module.java}) en tant que \texttt{Composant} composite : \texttt{public class Module extends Composant}, conform\'ement \`a la structure pr\'evue par le simulateur. Premi\`ere coquille fonctionnelle, l'impl\'ementation interne (\texttt{calculer()}) reste \`a \'etoffer.
  \item \textbf{Corrections cibl\'ees} pour fiabiliser la batterie de tests c\^ot\'e simulateur : champs \texttt{this.} explicit\'es dans \texttt{Lien.java} (\texttt{4c46433}), correction d'un retour de \texttt{ErreurIndex} et compl\'etude de \texttt{Composant.java} pour les tests JUnit (\texttt{9ef77d7}).
  \item \textbf{Mise \`a jour des scripts de build} (\texttt{LM/Build.sh}) : prise en compte des sources \texttt{tests projet long/simulateur} pour permettre \`a la fen\^etre de compiler avec ses nouvelles d\'ependances.
\end{itemize}

\subsection*{R\'esultats}

\begin{itemize}
  \item \textbf{Fen\^etre fonctionnelle de bout en bout} : 9~entr\'ees, 9~sorties, 4~portes AND c\^abl\'ees, propagation v\'erifi\'ee manuellement. C'est le premier livrable o\`u l'IHM ne ment plus sur l'\'etat du circuit.
  \item \textbf{Hi\'erarchie de composants \'etendue} : la classe \texttt{Module extends Composant} ouvre la voie \`a la composition de circuits hi\'erarchiques (un module r\'eutilisable comme un composant ordinaire), conform\'ement \`a ce que demande la grammaire SHDL c\^ot\'e parser.
  \item \textbf{Tests existants pass\'es au vert apr\`es modifications} : les ajustements sur \texttt{ErreurIndex} et \texttt{Composant} ont \'et\'e dict\'es par la batterie JUnit du simulateur, qui sert de garde-fou \`a chaque modification structurelle.
\end{itemize}

\subsection*{Difficult\'es}

La principale difficult\'e a \'et\'e d'absorber l'API du simulateur sans en \^etre l'auteur : comprendre la s\'emantique des trois \'etats (UP, DW, ND), savoir quand un \texttt{Lien} doit \^etre cr\'e\'e par le \texttt{Simulateur} plut\^ot qu'instanci\'e directement, quel composant d\'eclenche le calcul. La lecture des tests JUnit s'est r\'ev\'el\'ee plus rentable que celle des Javadocs incompl\`etes.

\section{It\'eration 3 --- int\'egration \`a la nouvelle API du simulateur}

\subsection*{Objectif}

Pendant l'it\'eration~2, l'axe simulateur a entrepris une \textbf{refonte profonde de son IR}~: \texttt{Lien~$\to$~Connecteur}, \texttt{TableauLien~$\to$~TableauConnecteur}, \texttt{FileListe~$\to$~FileSimulateur}, ajout d'un dictionnaire \texttt{DicoConnecteur} indexant les signaux par nom. Plus important pour mon axe, le concept de \og module composite \fg{} que j'avais commenc\'e \`a impl\'ementer (\texttt{Module extends Composant}) a \'et\'e remplac\'e par un trio \texttt{Structure}~/ \texttt{StructEntree}~/ \texttt{StructSortie}, plus expressif et coh\'erent avec la sortie du convertisseur de l'axe parser.

L'objectif de l'it\'eration~3 \'etait donc d'\textbf{adapter mes trois fichiers d'interface \`a cette nouvelle API} et de les int\'egrer dans le dossier \texttt{simulateur/} commun, plut\^ot que dans mon dossier personnel \texttt{LM/}.

\subsection*{Activit\'es r\'ealis\'ees}

\begin{itemize}
  \item \textbf{Migration vers le dossier \texttt{simulateur/}} : abandon du dossier personnel \texttt{LM/} et d\'eplacement de mes trois fichiers dans le module commun, ce qui les rend disponibles \`a l'ensemble de l'\'equipe et permet leur compilation conjointe avec les autres composants.
  \item \textbf{Adaptation \`a la nouvelle API du simulateur} (commit \texttt{97a810e}, +215~lignes) :
  \begin{itemize}
    \item r\'e\'ecriture de \texttt{FenetreSimulation.java} (151~lignes) pour utiliser \texttt{SimulateurCarresRonds} \`a la place de l'ancien \texttt{Simulateur} ;
    \item cr\'eation de \texttt{SimulateurCarresRonds.java} (53~lignes), couche d'adaptation entre l'IR du simulateur et la rang\'ee de boutons ronds/carr\'es. C'est la nouvelle fronti\`ere entre mon IHM et le moteur ;
    \item cr\'eation de \texttt{SimulationCarresRonds.java} (9~lignes) comme classe d'entr\'ee \texttt{extends Application}, qui d\'el\`egue \`a \texttt{FenetreSimulation}.
  \end{itemize}
  \item \textbf{Scripts de build et d'ex\'ecution} (commit \texttt{dbd3815}) : ajout de \texttt{simulateur/Build.sh} et \texttt{simulateur/Start.sh}. Le script de build compile l'ensemble des sources du simulateur (\texttt{And}, \texttt{Or}, \texttt{Not}, \texttt{Connecteur}, \texttt{DicoConnecteur}, \texttt{Multiplicateur}, \texttt{Structure}, \texttt{StructEntree}, \texttt{StructSortie}, \dots) plus mes trois fichiers d'interface, avec le \texttt{module-path} JavaFX correctement positionn\'e.
  \item \textbf{Maintenance de la documentation et du \texttt{.gitignore}} : exclusion du dossier de documentation g\'en\'er\'ee localement (\texttt{*.html}, \texttt{*.class}) pour \'eviter la pollution des branches.
  \item \textbf{Correction post-int\'egration} (10~mai) : les exemples \texttt{modules/and.shdl} et \texttt{modules/module.shdl} servant de tests de bout en bout pour l'application int\'egr\'ee \'etaient r\'edig\'es avec la cl\^oture \texttt{end} alors que la grammaire LL(1) finalis\'ee impose \texttt{end~module}. Correction faite, ce qui d\'ebloque le bouton \og simuler \fg{} de la fen\^etre principale pour ces deux exemples.
\end{itemize}

\subsection*{R\'esultats}

\begin{itemize}
  \item \textbf{Fen\^etre de simulation compatible avec la version courante du simulateur de l'\'equipe} : compilation et lancement reproductibles via \texttt{cd simulateur \&\& bash Build.sh \&\& bash Start.sh}.
  \item \textbf{Code regroup\'e dans \texttt{simulateur/}} : plus de duplication entre \texttt{LM/} et le dossier commun, plus de divergence d'API entre ma fen\^etre et le moteur de simulation.
  \item \textbf{Param\'etrage du circuit interne via une couche d'adaptation} (\texttt{SimulateurCarres\-Ronds}) : la fen\^etre n'a plus besoin de connaître la moindre classe interne du simulateur en dehors de cette interface, ce qui isole l'IHM des futurs refactorings du moteur.
\end{itemize}

\subsection*{Difficult\'es}

La principale difficult\'e a \'et\'e de naviguer une refonte de l'API qui ne m'avait pas \'et\'e annonc\'ee : plusieurs noms de classes avaient chang\'e, et mon esquisse de \texttt{Module} \'etait obsol\`ete. J'ai pris le parti de ne pas tenter de sauver mes propres classes mais de me caler sur la nouvelle API : c'est plus court \`a impl\'ementer, et cela \'evite que mon IHM devienne le seul module \`a faire diverger la conception. Cette d\'ecision a permis de livrer un \og adapter \fg{} (\texttt{SimulateurCarresRonds}) qui est aujourd'hui le seul point de couplage entre la fen\^etre et le moteur.

\section*{Bilan global du sprint 2}

\subsection*{R\'eussites}

\begin{itemize}
  \item Passage d'une maquette d\'econnect\'ee (sprint~1) \`a une fen\^etre de simulation \emph{r\'eelle\-ment fonctionnelle}, robuste \`a une refonte d'API entre les deux it\'erations du sprint.
  \item Fourniture de scripts \texttt{Build.sh} / \texttt{Start.sh} reproductibles, qui ont permis aux autres membres de l'\'equipe de v\'erifier ma fen\^etre sur leur poste sans \^etre familiers du \texttt{module-path} JavaFX.
  \item Capacit\'e d'adaptation aux d\'ecisions de l'\'equipe (renommages \texttt{Lien~$\to$~Connecteur}, abandon de mon \texttt{Module} au profit de \texttt{Structure}) sans cr\'eer de dette technique parall\`ele.
\end{itemize}

\subsection*{Limites et travaux pour la suite}

\begin{itemize}
  \item Le circuit affich\'e est toujours param\'etr\'e \og 9~entr\'ees, 9~sorties \fg{}, ind\'ependamment du module r\'eel charg\'e. Une it\'eration suivante devrait synchroniser ce dimensionnement avec le \texttt{Module} effectivement converti par le parser.
  \item La fen\^etre \texttt{FenetreSimulation} reste autonome (lan\c{c}able via \texttt{Start.sh}) et n'est pas encore appel\'ee depuis la \texttt{FenetrePrincipale} de la branche \texttt{GUI}, qui ouvre actuellement sa propre fen\^etre \texttt{FenetreSimulateur} via le bouton \og simuler \fg{}. Une convergence est \`a discuter pour \'eviter de maintenir deux fen\^etres concurrentes.
  \item La classe \texttt{Module} reste un squelette. Elle a \'et\'e d\'epass\'ee par les classes \texttt{Structure~/ StructEntree~/ StructSortie} de l'\'equipe ; soit elle est retir\'ee, soit elle devient un alias de la nouvelle structure.
\end{itemize}

\section*{Manuel utilisateur}

\textbf{Pr\'e-requis} : JDK 21+ (testé sur Temurin~25), JavaFX~17 fourni dans \texttt{javafx-sdk-17.0.18/} \`a la racine du d\'ep\^ot.

\textbf{Compiler et lancer la fen\^etre de simulation} (depuis la branche \texttt{simulation}) :
\begin{verbatim}
  cd simulateur
  bash Build.sh
  bash Start.sh
\end{verbatim}

\textbf{Utilisation} : la rang\'ee sup\'erieure de carr\'es repr\'esente les 9~entr\'ees ; cliquer pour basculer entre UP, DW et ND. La rang\'ee inf\'erieure de ronds affiche les sorties calcul\'ees apr\`es propagation \`a travers les composants enregistr\'es dans le \texttt{SimulateurCarresRonds}.

\end{document}
